\subsection{极小权}

命题 6. 设 \(\lambda \in \mathrm{P}\)，\(\mathcal{X}\) 是包含 \(\lambda\) 的 P 的最小 R-饱和子集。取如命题 5 中的范数 \(\| \cdot \|\)。以下条件等价：

\begin{enumerate}
\def\labelenumi{(\roman{enumi})}
\item
  \(\mathcal{X} = \mathrm{W}.\lambda\) ;
\item
  \(\mathcal{X}\) 的所有元素具有相同的范数；
\item
  对所有 \(\alpha \in \mathrm{R}\)，都有 \(\lambda(H_{\alpha}) \in \{0,1,-1\}\)。
\end{enumerate}

P 的每个非空 R-饱和子集都包含一个满足上述条件的元素 \(\lambda\)。

引入条件：

（ii'）对所有 \(\alpha \in \mathrm{R}\)，设 \(t\) 为介于 0 与 \(\lambda(H_{\alpha})\) 之间的任一整数，

\[
\| \lambda - t \alpha \| \geq \| \lambda \|.
\]

\begin{enumerate}
\def\labelenumi{(\roman{enumi})}
\tightlist
\item
  \(\Longrightarrow\) (ii) \(\Longrightarrow\) (ii')：显然。
\end{enumerate}

\((\mathrm{ii}^{\prime}) \Longrightarrow (\mathrm{iii})\)：假设条件 \((\mathrm{ii}^{\prime})\) 满足。设 \(\alpha \in \mathrm{R}\)。我们有 \(\| \lambda \| = \| \lambda - \lambda (H_{\alpha})\alpha \|\)，所以有 \(\| \lambda - t\alpha \| < \| \lambda \|\)，其中 \(t\) 严格位于 0 与 \(\lambda (H_{\alpha})\) 之间；因此不存在这样的整数，从而 \(|\lambda (H_{\alpha})| \leq 1\)。

\begin{enumerate}
\def\labelenumi{(\roman{enumi})}
\setcounter{enumi}{2}
\tightlist
\item
  \(\Longrightarrow\) (i)：假设条件 (iii) 满足。设 \(w \in W\) 且 \(\alpha \in R\)。则 \((w\lambda)(H_{\alpha}) = \lambda(H_{w^{-1}\alpha}) \in \{0,1,-1\}\)；因此，若 t 是介于 0 与 \((w\lambda)(H_{\alpha})\) 之间的整数，则 \(w\lambda - t\alpha\) 等于 \(w\lambda\) 或 \(s_{\alpha}(w\lambda)\)。这证明 W. \(\lambda\) 是 R-饱和的，所以 X = W. \(\lambda\)。
\end{enumerate}

设 Y 是 P 的非空 R-饱和子集。Y 中存在一个范数最小的元素 \(\lambda\)。显然 \(\lambda\) 满足条件 (ii')，从而得到命题的最后一个断言。

命题 7. 设 V 是有限维单 g-模，X 是 V 的权集，\(\lambda\) 是 X 的最高元（参见命题 5 (i)）。命题 6 的条件 (i)、(ii) 和 (iii) 等价于：

\begin{enumerate}
\def\labelenumi{(\roman{enumi})}
\setcounter{enumi}{3}
\tightlist
\item
  对所有 \(\alpha \in \mathrm{R}\) 以及所有 \(x \in \mathfrak{g}^{\alpha}\)，都有 \((x_{\mathrm{V}})^2 = 0\)。
\end{enumerate}

若这些条件满足，则 V 的所有权的重数都为 1。

若 (i) 满足，则 \(X = W \cdot \lambda\)，且所有权都与 \(\lambda\) 具有相同的重数（命题 2 的推论 2），换言之，重数为 1。此外，若 \(w \in W\) 且 \(\alpha \in R\)，则 \(w\lambda + t\alpha\) 是 V 的权必有 \(|t| \leq 1\)。因此，若 \(x \in g^{\alpha}\)，

\[
(x _ {\mathrm{V}}) ^ {2} (\mathrm{V} ^ {w (\lambda)}) \subset \mathrm{V} ^ {w (\lambda) + 2 \alpha} = 0,
\]

所以 \((x_{\mathrm{V}})^2 = 0\)，这证明了 (i) \(\Longrightarrow\) (iv)。

反过来，假设 (iv) 满足。设 \(\alpha \in \mathrm{R}\)，并给 V 赋予 \(\mathfrak{sl}(2,k)\)-模结构，定义此结构的元素为 \(X_{\alpha}, X_{-\alpha}, H_{\alpha}\)，这些元素属于 \(\mathfrak{g}\)。将 (iv) 应用于 \(x = X_{\alpha}\)，可知 \(\mathfrak{sl}(2,k)\)-模 V 的权属于 \(\{0,1,-1\}\)（参见 §1，第 2 节，命题 2 的推论）。特别地，\(\lambda(H_{\alpha}) \in \{0,1,-1\}\)，所以 (iv) \(\Longrightarrow\) (iii)。

命题 8. 假设 g 是单的。以 \(\alpha_{1},\ldots,\alpha_{l}\) 表示 B 的元素。设 \(\varpi_{1},\ldots,\varpi_{l}\) 是相应的基本权。设 \(H = n_{1}H_{\alpha_{1}} + \cdots + n_{l}H_{\alpha_{l}}\) 是 \(R^{\vee}\) 的最高根，J 是这样的 \(i \in \{1,\ldots,l\}\) 的集合：满足 \(n_{i} = 1\)。设 \(\lambda \in P_{++} - \{0\}\)。则命题 6 的条件 (i)、(ii) 和 (iii) 分别等价于以下每个条件：

\begin{enumerate}
\def\labelenumi{(\alph{enumi})}
\setcounter{enumi}{21}
\tightlist
\item
  \(\lambda(H)=1;\)
\end{enumerate}

\begin{enumerate}
\def\labelenumi{(\roman{enumi})}
\setcounter{enumi}{5}
\tightlist
\item
  存在 \(i \in \mathrm{J}\)，使得 \(\lambda = \varpi_i\)。
\end{enumerate}

\(\varpi_{i}\)（其中 \(i \in J\)）在 P(R) 中构成 P(R)/Q(R) 的非零元素的一个代表系。

设 \(\lambda = u_{1}\varpi_{1} + \cdots + u_{l}\varpi_{l}\)，其中 \(u_{1}, \ldots, u_{l}\) 是满足 \(\geq 0\) 且不全为零的整数。则 \(\lambda(H) = u_{1}n_{1} + \cdots + u_{l}n_{l}\)，且 \(n_{1} \geq 1, \ldots, n_{l} \geq 1\)，这立即给出 (v) 与 (vi) 的等价性。另一方面，\(\lambda(H) = \sup_{\alpha \in \mathrm{R}_{+}} \lambda(H_{\alpha})\)，且 \(\lambda(H) > 0\)，因为 \(\lambda\) 是 \(P_{++}\) 的非零元素。因此，条件 (v) 等价于条件 \(\lambda(H_{\alpha}) \in \{0, 1\}\) 对所有 \(\alpha \in R\) 成立，换言之，等价于命题 6 的条件 (iii)。

命题的最后一个断言由第 VI 章，§2，第 3 节，命题 6 的推论得出。

定义 1. 假设 \(\mathfrak{g}\) 是单的。\((\mathfrak{g},\mathfrak{h})\) 的极小权是 \(\mathrm{P}_{++} - \{0\}\) 中满足命题 6、7 和 8 的等价条件 (i)、(ii)、(iii)、(iv)、(v) 和 (vi) 的元素。

备注。假设 g 是单的。设 \(\Sigma^{\prime\vee}\) 是仿射 Weyl 群 \(\mathrm{W}_{a}(\mathrm{R}^{\vee})\) 的 Coxeter 图。回顾一下，\(\Sigma^{\prime\vee}\) 的顶点是 Coxeter 图 \(\Sigma^{\vee}\)（即 \(\mathrm{W}(\mathrm{R}^{\vee})\) 的 Coxeter 图）的顶点，再加上一个补充顶点 0。群 \(\mathrm{A}(\mathrm{R}^{\vee})\) 作用于 \(\Sigma^{\prime\vee}\)，并固定 0。群 \(\mathrm{Aut}(\Sigma^{\prime\vee})\) 典则同构于 \(\mathrm{A}(\mathrm{R}^{\vee})/\mathrm{W}(\mathrm{R}^{\vee})\) 与群 \(\Gamma_{C}\) 的半直积（参见第 VI 章，§2，第 3 节以及第 VI 章，§4，第 3 节）；显然 \((\operatorname{Aut} \Sigma'^{\vee})(0) = \Gamma_C(0)\)；并且 \(\Gamma_C(0)\) 由 0 以及 \(\Sigma^{\vee}\) 中与 \(\varpi_i\)（其中 \(i \in \mathbf{J}\)）对应的顶点组成（参见第 VI 章，§2，命题 5 以及第 3 节的备注 1）。总之，极小权是与 \(\Sigma^{\vee}\) 中可由 \(\operatorname{Aut}(\Sigma'^{\vee})\) 的某个元素作用于 0 而得到的顶点相对应的基本权。

采用第 VI 章图版 I 至 IX 的记号，由前述结果可推出，极小权如下：

对于 \(A_{l}\) 型（\(l \geq 1\)）：\(\varpi_{1}, \ldots, \varpi_{l}\)。

对于 \(B_{l}\) 型（\(l \geq 2\)）：\(\varpi_{l}\)。

对于 \(C_{l}\) 型（\(l \geq 2\)）：\(\varpi_{1}\)。

对于 \(D_{l}\) 型（\(l \geq 3\)）：\(\varpi_{1}, \varpi_{l-1}, \varpi_{l}\)。

对于 \(E_{6}\) 型：\(\varpi_{1}, \varpi_{6}\)。

对于 \(\mathrm{E}_7\) 型：\(\varpi_7\)。

对于 \(E_{8}, F_{4}, G_{2}\) 型，不存在极小权。

